\documentclass[11pt]{article}

\usepackage[margin=1in]{geometry}
\usepackage{microtype}
\usepackage{booktabs}
\usepackage{hyperref}
\hypersetup{
    colorlinks=true,
    urlcolor=blue!60!black,
    linkcolor=black,
    citecolor=black
}
\usepackage{parskip}
\usepackage{enumitem}
\setlist{itemsep=2pt,topsep=2pt}

\title{Sub-Decision Prosodic Sensitivity in Speech-LLMs\\
\large Problem Statement and Project Plan}
\author{Piyush Kumar Agarwal}
\date{}

\begin{document}
\maketitle

\section*{The Research Question in One Sentence}

When a speech-capable large language model gets a contrastive-stress
question \emph{wrong}, does the acoustic stress in the speaker's voice
still shift the model's internal answer distribution --- even if it
doesn't shift the final answer?

\section*{Why This Question Matters}

Speech-LLMs are increasingly deployed in voice assistants, transcription
tools, and accessibility software. Prior work (StressTest, HUJI 2025)
shows these models often pick the wrong answer when meaning hinges on
prosody. But a \emph{wrong final answer} is a coarse signal. The model
might be:

\begin{itemize}
    \item \textbf{Truly deaf to prosody} --- its internal probabilities
    don't move at all.
    \item \textbf{Partially sensitive} --- its probabilities shift in
    the correct direction, but not enough to change the top choice.
\end{itemize}

These two failure modes call for very different engineering fixes. The
first implicates the audio encoder (``the ears''); the second
implicates the language model backbone (``the brain''). This project
distinguishes them by reading the model's hidden probability
distribution instead of just its top-1 answer.

\section*{Scope and Constraints}

\begin{itemize}
    \item \textbf{Inference only} --- no training, no fine-tuning.
    \item \textbf{Primary model:} Qwen2-Audio-7B-Instruct, loaded in
    4-bit quantization.
    \item \textbf{Hardware:} single 16~GB GPU (Google Colab T4).
    \item \textbf{Timeline:} $\sim$6 weeks, part-time.
    \item \textbf{Benchmark:} StressTest (\texttt{slprl/StressTest} on
    HuggingFace) --- 218 audio items, each a sentence whose meaning
    depends on which word the speaker stresses.
\end{itemize}

\section{Phase 1 --- Foundation: Setting the Baselines}

\textbf{Goal.} Build a controlled environment so any signal we
attribute to stress is defensibly \emph{not} an artifact of noise,
text, or pipeline bugs.

\subsection*{Step 1.1 --- Verify the Benchmark}

Confirm the StressTest dataset loads, has the expected schema (audio
at 16~kHz, two possible answers per item, a ground-truth label), and a
usable license (CC-BY-NC-4.0 for academic work).

\textbf{Rationale.} Every downstream measurement assumes the dataset
is what we think it is. A two-hour verification step at the start
prevents weeks of analysis built on misread columns.

\subsection*{Step 1.2 --- Build the Acoustic Noise Floor}

Apply small, \textbf{meaning-preserving} perturbations to each audio
clip --- a mild gain change, low-level additive noise, a $\pm 5\%$
time-stretch --- and measure how much the model's internal probability
for the correct answer shifts purely from those perturbations.

\textbf{Rationale.} Any model's output distribution wobbles slightly
when the input changes, even when the change is semantically
meaningless. We need to know how big that ``baseline wobble'' is
before we can claim a stress-induced shift is real. The
pre-registered rule: a stress shift must exceed the \textbf{95th
percentile} of the noise-floor distribution to count as meaningful.

\textit{Design note.} An earlier plan considered using same-speaker
duplicate takes from the benchmark as the noise floor. On inspection,
StressTest's multiple versions of each sentence carry \emph{different}
stress patterns --- they are the experimental stimuli, not duplicate
takes. So we use synthetic perturbations instead, which are also
fully reproducible.

\subsection*{Step 1.3 --- Text-Only and Whisper Falsification Checks}

Feed the written \textbf{transcript} of each audio clip into:

\begin{enumerate}
    \item A text-only LLM (no audio modality at all).
    \item A cascade pipeline: Whisper transcribes the audio
    $\rightarrow$ text-LLM answers from the transcript.
\end{enumerate}

\textbf{Rationale.} If either of these systems can solve StressTest
from text alone, then audio prosody is \emph{not} the only path to the
answer, and the experiment is ill-posed. Both pipelines should
perform near chance. This is a falsification check on the entire
study premise.

\section{Phase 2 --- Core Experiment: Logit Extraction}

\textbf{Goal.} Measure not just what the model \emph{says}, but the
probability distribution behind what it says.

\subsection*{Step 2.1 --- Input the Audio Pairs}

Each StressTest item is part of a pair: same words, different
stressed word, opposite correct answer. Feed both audio versions
through the model.

\subsection*{Step 2.2 --- Strict Forced-Choice Prompting}

Ask the model a single forced-choice question with two answer
options, and prompt it so the next predicted token is constrained to
one of two specific tokens. We use two answer formats:

\begin{itemize}
    \item \textbf{Format A:} options labeled \texttt{A.} and
    \texttt{B.}; next token is \texttt{\ A} or \texttt{\ B}.
    \item \textbf{Format B:} options labeled \texttt{1.} and
    \texttt{2.}; next token is \texttt{1} or \texttt{2}.
\end{itemize}

The two formats let us test whether any observed effect is genuine
prosodic sensitivity or an artifact of how the options were labeled.

\subsection*{Step 2.3 --- Stop and Extract the Logits}

Instead of letting the model generate a full sentence, we halt after
the \textbf{very first token} and read the raw logit values for
exactly the two competing answer tokens. A two-way softmax over those
two logits gives a clean probability for each candidate answer.

\textbf{Rationale.} A free-form generated answer collapses a rich
probability distribution into a single argmax decision. Reading the
logits directly preserves the sub-decision information that is the
whole point of this study. Restricting the softmax to the two answer
tokens (rather than the full vocabulary) gives a well-defined
``probability of A given the question is A or B.''

\textit{Pipeline integrity check.} Independently run
\texttt{model.generate()} for one token on a sample of items and
assert that its output token matches the argmax of our extracted
logits. If they disagree, the extraction code has a bug and
everything stops until it's fixed.

\section{Phase 3 --- Safety Checks and Branching Results}

\textbf{Goal.} Decide, per item, whether the model exhibited a real
reaction to stress.

\subsection*{Step 3.1 --- Calculate the Shift}

For each pair, compute how much the probability of the \emph{correct}
answer changed between the two audio versions. This is the
``stress-induced shift.''

\subsection*{Step 3.2 --- Filter Against the Noise Floor (Branching Point)}

Compare the stress-induced shift to the noise-floor distribution from
Step 1.2.

\begin{itemize}
    \item \textbf{Path A --- Significant Shift.} The shift exceeds the
    95th-percentile noise threshold. The model \emph{noticed} the
    stress at the distributional level. If it still picked the wrong
    final answer, this is a \textbf{measurable-but-decision-
    insufficient} influence: the signal was there but underweighted.
    \item \textbf{Path B --- Negligible Shift.} The shift falls within
    the noise floor. We categorize the model as showing \textbf{no
    measurable influence} for this item.
\end{itemize}

\subsection*{Step 3.3 --- Format Swap for Robustness}

Re-run the entire experiment with Format B (1/2) and check that
whatever directional effect we saw in Format A (A/B) survives. A real
prosodic effect should not depend on whether the answer was called
``A'' or ``1.''

\section{Phase 4 --- Diagnostic Probe: Pinpointing the Failure}

\textbf{Goal.} When the model fails (especially Path~B items),
determine \emph{where} in the architecture the failure lives.

\subsection*{Step 4.1 --- Text-Rescue Test}

Re-present each failed item to the model, but this time encode the
stress \textbf{in the transcript} by capitalizing the stressed word
in ALL CAPS (e.g., \texttt{"I NEVER said he stole it."}).

\begin{itemize}
    \item \textbf{Outcome A --- The ``Ears'' Failed.} The model now
    answers correctly. The language backbone can use prosodic
    information when it arrives in text form; the bottleneck is the
    upstream audio pathway (encoder, projector, fusion).
    \item \textbf{Outcome B --- The ``Brain'' Failed.} The model
    still gets it wrong even with the explicit text hint. The
    limitation lives downstream of the audio pathway, in how the
    language model integrates the stress signal into its reasoning.
\end{itemize}

\textbf{Rationale.} This is the diagnostic payoff of the whole study.
It converts an opaque ``the model got it wrong'' into an actionable
architectural attribution.

\section{Phase 5 --- Final Verdict: Statistical Proof and Conclusion}

\textbf{Goal.} Convert per-item measurements into a defensible
aggregate claim.

\subsection*{Step 5.1 --- Statistical Testing}

\begin{itemize}
    \item \textbf{Wilcoxon signed-rank test} on paired shifts
    (matched pairs of audio stimuli). Chosen because it's
    non-parametric and appropriate for paired data without
    distributional assumptions.
    \item \textbf{Bootstrap 95\% confidence intervals} on the shift
    magnitude and the decision-flip rate, to characterize effect
    size and uncertainty.
\end{itemize}

\textit{Honesty note.} With $\sim$109 pairs, statistical power is
modest. Confidence-interval \emph{widths} will be reported as
prominently as $p$-values, so the conclusion is calibrated to the
data we actually have.

\subsection*{Step 5.2 --- Final Graded Conclusion}

For the primary model (Qwen2-Audio-7B), pair the quantitative finding
from Step~5.1 with the architectural diagnosis from Phase~4 to produce
a per-item characterization of the model's prosodic sensitivity.

\section*{Headline Outputs of the Study}

\begin{enumerate}
    \item \textbf{Decision-flip rate} --- fraction of pairs where the
    model's top answer flips in the correct direction when stress
    flips.
    \item \textbf{Stress-induced shift magnitude vs.\ acoustic noise
    floor} --- the distribution of sub-decision shifts, plotted
    against the noise threshold.
    \item \textbf{Architectural attribution} --- per-item
    categorization of ``ears failed'' vs.\ ``brain failed'' from the
    text-rescue probe.
\end{enumerate}

\section*{What This Study Will \emph{Not} Claim}

\begin{itemize}
    \item It will not claim Qwen2-Audio is universally
    prosody-blind or prosody-aware. The claim is specific to
    \textbf{contrastive lexical stress} as operationalized by
    StressTest.
    \item It will not claim training-time fixes. The project is
    inference-only and diagnostic.
    \item It will not generalize beyond English or beyond the
    speakers in StressTest.
\end{itemize}

\end{document}

